% Source: https://artofproblemsolving.com/community/c3h1405026_volume_traced
% Stable identity: aops-topic-1405026-q01
Let $P$ and $Q$ be the fixed vertices, so the axis is the diagonal $PQ$, and let $X$ and $Y$ be the other two vertices. The square is symmetric about the axis: it consists of triangles $PQX$ and $PQY$. Under a half turn, each point at distance $d$ from the axis sweeps a semicircle of radius $d$, and the two halves together sweep the full circle. So the solid is the solid of revolution of triangle $PQX$ about $PQ$:

\begin{asy}
usepackage("amsmath");
size(4.5cm);
defaultpen(fontsize(10pt));
real s=1, r=s/sqrt(2);
pair L=(-r,0),R=(r,0),Tp=(0,r),Bt=(0,-r);
draw(L--Tp--R--Bt--cycle,linewidth(0.9));
draw(xscale(0.36*r)*yscale(r)*arc((0,0),1,-90,90),blue);
draw(xscale(0.36*r)*yscale(r)*arc((0,0),1,90,270),blue+dashed);
draw(L--R,gray+dashed);
dot(L);dot(R);dot(Tp);dot(Bt);
label("$P$",L,W); label("$Q$",R,E); label("$X$",Tp,N); label("$Y$",Bt,S);
label("$s$",(L+Tp)/2,NW);
draw((0,0)--(0,r),red); label("$r$",(0,0.45*r),W,red+fontsize(9pt));
label("\mbox{}",(min(currentpicture,user=true)+max(currentpicture,user=true))/2,black); // keeps the text after the figure black when LaTeX typesets the labels inline
\end{asy}

it is a \textbf{double cone}, two congruent cones with apexes $P$ and $Q$ glued along a common circular base through $X$ and $Y$. The base radius $r$ (red in the figure) is the distance from $X$ to the axis, half the diagonal, $r=\dfrac{s}{\sqrt2}$, and each cone has the same height $\dfrac{s}{\sqrt2}$. So the volume is
\[
2\cdot\dfrac{1}{3}\pi\left(\dfrac{s}{\sqrt2}\right)^2\cdot\dfrac{s}{\sqrt2}=\dfrac{2\pi s^3}{3\cdot2\sqrt2}=\boxed{\dfrac{\sqrt2\,\pi s^3}{6}} .
\]
